% !TeX root = ../papers/nc-p3.tex
\section{平面二次曲線成分の一次・二次生成元による全次数判定}
\label{sec:conic-component-as-test}
同じ点ファイバーの仮定で、全位置で平面二次曲線になる純な部分スキーム $Z\subset X$ を印とする。
四移動によるその軌道 $\mathcal C$ が有限とする。特に滑らかな二次曲線成分はこの条件を満たす。
$D(Z)$ は全切断商、$I(Z)$ はその核である。平面式を $x_i(Z)\in V_i$ とし、スカラー倍は任意とする。

\begin{theorem}\label{thm:conic-component-as-test}
各 $Z\in\mathcal C$ と位置 $i$ で
$y_i(Z)\in I_{i,i+2}(Z)\setminus x_i(Z)V_{i+1}$ を選ぶ。
標準Koszul左右AS正則性は次の有限条件と同値であり、どの持上げ $y_i$ を使ってもよい。
軌道の置換 $\pi$ が存在し、全 $i,Z$ について $Z'=\pi Z$ と書くと
\[
 x_i(Z) a + y_i(Z) b=0\quad\text{in } A_{i,i+3},\qquad
 0\ne b\in k x_{i+2}(Z'),\quad
 a\in I_{i+1,i+3}(Z')\setminus V_{i+1}b
\]
を満たす係数がある。検査する位置は $4|\mathcal C|$ 個で十分である。
\end{theorem}
\begin{proof}
二次曲線の各偏極はその平面の $\mathcal O(1)$ である。滑らかなら $\mathbb P^1$ の $\mathcal O(2)$、
二直線なら各直線上次数1で木の貼合せ、二重直線なら直線上の冪零部分 $\mathcal O(-1)$ の
$H^1$ が零なので、いずれも全位置で同じ束 $L$ となる。
平面での一つの二次式の解消から
$\omega_Z=L^{-1}$、$h^0(L^n)=2n+1$ ($n\ge0$) で、切断環は一次生成である。
$D(Z)$ はこの環の反復となる。$I_{i,i+2}$ の次元は5である。

直線の場合と同じ局所分裂束複体を用いる。$n\ge4$ なら全項の $H^1$ は零。
$n=3$ では項4に $H^1(L^{-1})=k$ があり、項2の切断ホモロジーを一元与える。
$n=2$ では項4から項3への $H^1$ の射が三次元から四次元への単射である。
その双対は四文字から三つの平面切断への全射である。従って項1の切断ホモロジーだけが一次元残る。
$n=1$ では平面式一本が項1に残る。$n=0$ は単純のみである。
よって右の第一Torは次数1,2に一元ずつ、第二Torは次数3の一元の商であり、反対側も同じ。
特に $x_i V_{i+1}$ は四次元で、右イデアルは $x_i,y_i$ で生成され、関係は次数3に高々一つ。
左イデアルも終点に応じて同様である。

標準ASならこれが長さ2の極小解消となる。最終列を $(a,b)^t$ として双対化する。
定理\ref{thm:elliptic-quartic-as-test}の局所双対は軌道を折り返した同じ証明で適用でき、
第二Extの最小生成元は左から四移動で運ばれた印 $Z'$ のイデアルで零化される。
正Torの遷移消滅には $[N,N+2]$ の内部次数範囲を使えばよい。
従って次数1において $x_{i+2}(Z')$ は $b$ の倍数、次数2では
$V_{i+1}b+k a$ が五次元の $I_{i+1,i+3}(Z')$ を含む。以上から表示の全条件を得る。
置換はねじれまたはその逆が全軌道に誘導するものを取る。
$y_i$ を非零倍と $x_i c$ の加算で取り替えても、$a$ の変化は倍率と $V_{i+1}b$ の元だけである。

逆に有限条件を仮定する。$b\ne0$ より表示の関係は上で得た高々一つの最小関係を実現する。
固定終点で $h_n=\dim A_{j-n,j}$、$d_n=2n+1$ ($n\ge0$、負なら零) とおく。
右表示のEuler式は
\[
 h_n-h_{n-1}-h_{n-2}+h_{n-3}=d_n+e_n,\qquad e_n\ge0 ,
\]
ここで $e_n$ は末項の核次元である。一方 $m_n(Z)=\dim(A_{i,j}/x_i(Z)A_{i+1,j})$、$i=j-n$ とする。
この商から $D(Z)_{i,j}$ への核は $y_i(Z)A_{i+2,j}$ の像であり、
表示の関係により $y_i(Z)b A_{i+3,j}$ は消える。従って
$m_n(Z)\le d_n+m_{n-2}(\pi Z)$ である。$h_n\le h_{n-1}+m_n(Z)$ と合わせると上界は
\[
 \frac{H_D}{(1-t)(1-t^2)}=(1-t)^{-4}.
\]
先のEuler式に同じ非負係数の逆多項式 $1/((1-t)(1-t^2))$ を掛ければ下界もこれに等しい。
よって $e_n=0$、各不等式で等号、$x_i(Z)$ の左乗法は単射である。
反対側でも同じ上界の議論が通る。実際、終点 $i+3$ で $a,b$ は $Z'$ の左イデアルの二次・一次の
最小生成元であり、$x_i(Z) a=-y_i(Z)b$ を使い、置換を逆向きに辿ればよい。
既知の $H_A$ はこの上界と等しいから、全ての $x_i(Z)$ の右乗法も単射である。

右の長さ2の自由解消を双対化すると、最初の写像は $x_i$ の右乗法の単射性から単射で、
最後の余核は $a,b$ が生成する左イデアルによる商、すなわち $D(Z')$ の終点部分である。
中央のホモロジーも先のEuler等式（内部次数を末項に合わせる）で消える。
従って右切断商加群 $e_iD(Z)$ からの最初の二Extは零、反対側も同じ。
既出の環変更・Frobenius双対・Euler式で標準Koszul左右AS正則性が従う。
\end{proof}

\begin{corollary}\label{cor:conic-component-atlas}
滑らかな二次曲線成分を含む全窓は有限条件付き二窓アトラスを持つ。
同じ三つの座標平面に $[s^2:st:t^2:0]$ を取り、二辺で対角の二次曲線を消す
次元11の双一次式空間から六平面を選ぶ。
系\ref{cor:line-component-atlas}の二窓復元・点ファイバー・次元の検査と本定理を用いる。
指定を忘れた同型は第二窓の各二次曲線成分を順に標準化し、最初の三位置で
曲線上同じ $\mathrm{PGL}_2$ を誘導する平面安定化群の三変換により
二辺が同時に移るかで判定できる。一移動・二移動と同じ比較をすれば最小周期が決まる。
\end{corollary}
\begin{proof}
二平面の二次曲線切断の積は $H^0(\mathcal O_{\mathbb P^1}(4))$ 全体を張る。
六平面の周囲の次元は従って $16-5=11$ である。全ての曲線成分は位置ごとの次数を保ち、
二次曲線上の写像はその完全二次線型系により平面の射影変換に延びる。
残りは二窓の延長と有限個の成分の列挙による。
\end{proof}

\begin{corollary}\label{cor:full-quintic-point-window-atlas}
点ファイバーを持つ五次点グラフの全標準AS窓は、非被約構造と埋込み点を含め、
次の互いに素な三行で被覆される。
\[
\begin{array}{c|l}
\text{被約台の成分}&\text{全次数AS性と添字保持同型}\\ \hline
\text{直線あり}&\text{系}\ \ref{cor:line-component-atlas}\\
\text{直線なし、二次曲線あり}&\text{系}\ \ref{cor:conic-component-atlas}\\
\text{両方なし}&\text{滑らかな楕円五次の二類、定理}\ \ref{thm:elliptic-quintic-regularity}
\end{array}
\]
最初の二行では純部分の次数5を課す。これは射影像の一次のHilbert多項式の最高係数で検査できる。
種数二の可約五次での長さ二の有限構造も個別の接型を指定せずこの判定に入る。
さらに任意の次数について、直線・二次曲線も整な純部分もない残りの曲線台は、
三次二成分、三次と四次の二成分、または一本の整三次上の二重構造だけである。
\end{corollary}
\begin{proof}
各成分の一般長と次数の積の和が片側次数である。次数1の整曲線は直線、次数2は滑らかな平面二次曲線である。
これらがなければ二成分以上または一般長2以上には次数6以上を要し、
総次数上限7の下では記載の三通りに尽きる。一成分で一般長1なら純部分の冪零イデアルは
有限台なので零であり、系\ref{cor:integral-hull-high-degrees}を適用できる。
最初の二行では指定を忘れる際、同一窓の複数の印も同型判定で同一視される。
\end{proof}
